\section{Introduction}
A photograph of a facial skin condition contains more than a lesion. Its appearance depends on illumination, camera distance, skin texture, treatment history, facial location and the source from which the image was obtained. These differences matter to machine learning because a classifier has no inherent reason to attend to clinically meaningful features. When one disease category is collected mainly from one website or photographic setting, a model may recognize the source rather than the condition. An improvement on a random image split can therefore be misleading, particularly if resized versions of the same photograph appear in different partitions.

Clinical dermatology datasets face persistent problems of representation and source variation. The Fitzpatrick17k work drew attention to skin-type coverage and classification differences across skin types \citep{groh2021fitzpatrick}. SCIN adopted a different acquisition strategy, collecting user-contributed photographs of everyday skin concerns \citep{ward2024scin}. These efforts make valuable images available, but neither automatically supplies a diagnosis-verified, patient-independent six-class collection of facial photographs. A labeled image can depict a body site other than the face, more than one skin finding, or a diagnosis incompatible with the category to which it has been assigned.

Generative models offer a possible way to enrich the visual range of a limited dataset. Latent diffusion models synthesize images by denoising compressed representations under text conditioning \citep{rombach2022latent}. Yet a realistic generated face is not necessarily an informative dermatology sample. The prompt \emph{acne} leaves lesion count, size, inflammation and anatomical distribution unspecified. A more descriptive prompt can request those features, but it may also exceed the text encoder's effective conditioning capacity or produce incorrect lesions. Whether detailed prompts help a classifier is therefore a question for controlled evaluation, not an assumption about the generator's realism.

This paper presents \emph{DermSynReal}, an experimental framework designed around one central question: \emph{does morphology-guided synthetic-image augmentation change performance on unseen real facial photographs, beyond what is achieved with real-only learning or generic disease-name synthetic prompts?} We distinguish three effects. RQ1 concerns augmentation itself, comparing real-only training with real-plus-morphology training. RQ2 concerns prompt design, comparing generic and morphology-guided samples when the generator, sampling budget and learning schedule are matched. RQ3 concerns selection, testing whether quality- or diversity-based choice improves over a random matched-count sample from the same morphology pool. An optional synthetic-only arm tests the real-versus-synthetic domain gap.

The contribution is a reproducible \emph{evaluation design} accompanied by implementation rather than a claim of superior classification. It brings together provenance checks, duplicate-aware real-image grouping, a structured prompt library, base diffusion and optional training-only LoRA adaptation, five visual encoders, and paired evaluation on a frozen real test. This separation matters: a benefit attributed to clinical prompt detail is confounded if the generic and detailed images come from different generator checkpoints. Similarly, a selection gain cannot be isolated if one arm contains more training examples or was tuned on the real test.

The scope is a six-output photographic research classifier, not medical diagnosis, treatment guidance or prospective clinical validation. The intended labels contain clinically important subtypes, especially within fungal, pigmentary and malignant or premalignant conditions. Final conclusions must acknowledge uncertain diagnoses, image licenses, unknown patient identities and the limits of a single collected test domain. Numerical conclusions are left open until those questions have been addressed through actual runs and review.
